\section{ORDER BY dan LIMIT}

Klausa \code{ORDER BY} mengurutkan hasil query berdasarkan satu atau lebih kolom. Secara default, pengurutan bersifat \code{ASC} (ascending/naik); tambahkan \code{DESC} untuk urutan menurun. Anda dapat mengurutkan berdasarkan beberapa kolom: \code{ORDER BY kolom1 ASC, kolom2 DESC} --- MySQL akan mengurutkan berdasarkan kolom pertama terlebih dahulu, lalu kolom kedua untuk baris yang bernilai sama di kolom pertama \cite{mysqlDocs}.

Klausa \code{LIMIT} membatasi jumlah baris yang dikembalikan. Ini sangat berguna untuk \textbf{pagination} (menampilkan data per halaman) di aplikasi web. Sintaks: \code{LIMIT jumlah} untuk membatasi jumlah baris, atau \code{LIMIT offset, jumlah} untuk melewati sejumlah baris pertama. Misalnya, halaman 2 dengan 10 item per halaman: \code{LIMIT 10, 10} (lewati 10, ambil 10). Kombinasi ORDER BY dan LIMIT adalah pola yang sangat umum dalam query aplikasi web \cite{mysqlGettingStarted}.

\begin{webcode}[language=SQL, caption={ORDER BY dan LIMIT untuk pengurutan dan pagination}]
-- Urutkan berdasarkan harga (termurah dulu)
SELECT nama, harga FROM produk ORDER BY harga ASC;

-- Urutkan berdasarkan tanggal terbaru
SELECT * FROM produk ORDER BY created_at DESC;

-- Top 5 produk termahal
SELECT nama, harga FROM produk ORDER BY harga DESC LIMIT 5;

-- Pagination: halaman 1 (10 item)
SELECT * FROM produk ORDER BY id LIMIT 0, 10;

-- Pagination: halaman 2 (10 item)
SELECT * FROM produk ORDER BY id LIMIT 10, 10;

-- Pagination: halaman 3 (10 item)
SELECT * FROM produk ORDER BY id LIMIT 20, 10;
\end{webcode}

